% =============================================================================
% SECTION 9: DISCUSSION
% =============================================================================

\section{Discussion}
\label{sec:discussion}

The central finding of this work is that constraint tightness and answer accuracy are
independently controlled variables. Tighter constraints reduce the search space but do
not improve Hits@1. This result is counterintuitive: reducing the number of irrelevant
paths should make the LLM's task easier. Calibration and the v2 ablation isolate why it
does not: the gates themselves are the source of the accuracy cost, and step-wise trie
rebuilding is an architectural risk that costs 25 points of Hits@1, more than the 4.0--4.7 points the gates cost.

\subsection{Interpreting the Non-Monotone Result}
\label{sec:discussion:interpretation}

The relationship between SIR and Hits@1 is non-monotone. There exists an optimal
constraint level that balances diversity and relevance. Below this level (too loose), the
beam search is overwhelmed by irrelevant paths. Above this level (too tight), the beam
search loses diversity. The optimal point depends on the LLM's capacity to discriminate
relevant from irrelevant structure, the ontology's coverage, and the question's
complexity.

For practitioners, the TypeOracle provides a way to measure where a system sits on this
spectrum. SIR quantifies constraint tightness independently of accuracy. A system with
SIR close to 1.0 is under-constrained; a system with SIR close to 0 is over-constrained.
The current results show that even the modest SIR reduction of DCA-Trie v1 (from 1.0 to
0.138, a 13.8\% path reduction) costs 4.0pp of Hits@1, so the operating point that
balances diversity and relevance must be found empirically for each task rather than
assumed from the oracle's selectivity.

For researchers, the non-monotone result means that Hits@1 alone is insufficient to
evaluate constrained decoding systems. A system that achieves high Hits@1 with high SIR
may be over-constrained and fragile: small changes to the oracle could break it. SIR
provides the missing diagnostic. Reporting both metrics reveals whether a system's
accuracy depends on structural constraint (fragile) or on the LLM's reasoning capacity
(robust).

\subsection{Comparison with Prior Work}
\label{sec:discussion:comparison}

GCR~\citep{luo2025-graph-constrained-reasoning} achieves 92.6\% Hits@1 on WebQSP with
SIR $= 1.0$ using a two-stage pipeline; our single-stage reproduction reaches 91.6\%
on the full test set, on par with DoG's published 90.99\%, and the residual gap to the
published figure is the ChatGPT second stage. GCR's strength is that it preserves full
trie diversity: the beam search sees all structurally valid paths, and the LLM learns to
discriminate relevant from irrelevant structure during fine-tuning. DCA-Trie's oracle
removes this diversity, and the LLM cannot compensate.

DoG~\citep{li2024-decoding-on-graphs} constrains decoding to well-formed chains and
demonstrates that structural validity improves faithfulness. DoG does not measure SIR,
but its chains are structurally valid by construction (SIR $= 1.0$ for valid chains).
Our v2 result is directly relevant to DoG: v2 follows DoG's step-wise pattern of
expanding the trie after each entity commitment, and that architecture alone costs 25
points of Hits@1, independent of the gates. The step-wise trie-rebuilding design, shared
with DoG, is the risk our controlled experiments identify.

ToG~\citep{sun2024tog} uses an LLM-based oracle to explore the KG and generate
reasoning paths. ToG's oracle is adaptive (it reasons at each step), similar in spirit
to DCA-Trie v2. Our v2 results show that per-entity oracle evaluation compounds error,
driving Hits@1 25 points below the baseline. The lazy v3 variant achieves the adaptivity
that motivates such step-wise methods, but in a single stable decoding pass, at 3.5$\times$
lower construction cost than the static trie.

GNN-RAG~\citep{mavromatis2022rearev} retrieves relevant subgraphs using GNNs and
achieves 90.7\% Hits@1 on WebQSP. GNN-RAG's retrieval step is analogous to our oracle:
both filter the KG before decoding. The key difference is that GNN-RAG uses a learned
retriever (with implicit soft constraints) while DCA-Trie uses symbolic gates (hard
constraints). Our results suggest that soft constraints may preserve more diversity than
hard constraints, which is consistent with GNN-RAG's strong performance.

\subsection{Limitations}
\label{sec:discussion:limitations}

The paired comparison runs on a controlled 300-question WebQSP subset (seed 42) to keep
every condition matched on identical inputs; the full test set is used for the
reproduction study and gate decomposition. A larger paired run would narrow the
confidence intervals on the significance tests. The published GCR and DoG numbers are
reported as reference, with the protocol differences (two-stage versus single-stage,
$L{=}2$ index versus full set) stated explicitly; DCA-Trie is not claimed to exceed
them.

All experiments use Llama-3.1-8B fine-tuned on GCR data. Larger models (Llama-70B,
GPT-4) may handle diversity loss better, potentially reducing the Hits@1 gap. Smaller
models (Llama-7B without fine-tuning) may be more sensitive to diversity loss. The
generalizability across model scales is untested.

The type gate's regex classifier is the weakest component. A neural type classifier
would improve FNR and potentially narrow the Hits@1 gap. However, the calibration
result shows the full gated-number gap is attributable to the gates, so improving the
classifier, not the decoding mechanism, is the lever.

Freebase's ontology is fixed. Dynamic KGs (e.g., Wikidata with frequent updates) would
require the TypeOracle to update its type and range metadata. The current system does not
support incremental ontology updates.

The LLM is fine-tuned on GCR's training data, which uses a static trie. Fine-tuning on
DCA-Trie's filtered trie may teach the model to handle reduced diversity, potentially
narrowing the Hits@1 gap, the most promising direction for future work.
